\section*{Chapter \ref{ch:b1:is-lm} --- The IS--LM Model}
% Numbers recomputed in code/bachelor-1/24-is-lm/exercises/solutions.py (tested).

\begin{solution}{exo:b1:is-lm:is-ex}
$Y = 50 + 0.75(Y - 100) + 300 - 25r + 100$, so $0.25Y = 375 - 25r$ and $Y = 1500 - 100r$.
\end{solution}

\begin{solution}{exo:b1:is-lm:lm-ex}
$500 = Y - 100r$: $Y = 500 + 100r$, or $r = (Y - 500)/100$.
\end{solution}

\begin{solution}{exo:b1:is-lm:read}
In 2015 both were at or near the bound; Germany's rate was slightly
negative. The bound is not exactly zero: holding currency has costs
(storage, safety), so rates can go a little below zero, but not far.
\end{solution}

\begin{solution}{exo:b1:is-lm:sign}
With LM, $r$ rises; with a horizontal MP, $r$ does not change unless the bank
moves it.
\end{solution}

\begin{solution}{exo:b1:is-lm:eq}
$1500 - 100r = 500 + 100r$: $r = 5$ and $Y = 1000$. A rise $\Delta G$ shifts IS to
$Y = 1500 + 4\Delta G - 100r$, so $r = 5 + 0.02\Delta G$ and $Y = 1000 + 2\Delta G$: the multiplier is 2,
half the Keynesian-cross multiplier of 4.
\end{solution}

\begin{solution}{exo:b1:is-lm:monetary}
LM becomes $Y = 500 + \Delta M + 100r$; then $r = 5 - \Delta M/200$ and $Y = 1000 + \Delta M/2$: $0.5$.
\end{solution}

\begin{solution}{exo:b1:is-lm:trap-ex}
Ignoring the floor, $400 - 100r = 500 + 100r$ gives $r = -0.5$; with the floor $r = 0$ and
$Y = 400$. Doubling $M$ moves the \omterm{def:b1:is-lm:lm}{LM curve} right but its flat part stays at
zero: $Y$ stays 400.
\end{solution}

\begin{solution}{exo:b1:is-lm:mp-ex}
IS becomes $Y = 1340 - 100r$; $Y = 1000$ requires $r = 3.4$: a cut of $1.6$ points.
\end{solution}

\begin{solution}{exo:b1:is-lm:assets}
$2.239/0.891 = 2.51$ in 2008 and $4.487/0.891 = 5.04$ by 2015. With the rate at its
floor, LM is flat: more money leaves $r$ (and $Y$) where they are.
\end{solution}

\begin{solution}{exo:b1:is-lm:general}
As in the proof of \cref{prop:b1:is-lm:multipliers}. From $L_Y dY + L_r dr = 0$,
$dr/dG = -(L_Y/L_r)\,dY/dG > 0$.
\end{solution}

\begin{solution}{exo:b1:is-lm:crowding}
$dY/dG = 1/(1 - C' + I'L_Y/L_r)$. If $L_r \to -\infty$ (flat LM, a trap) or $I' \to 0$ (investment
insensitive to $r$), the term vanishes and $dY/dG = 1/(1 - C')$: no \omterm{def:b1:long-run-model:crowding-out}{crowding out}. If
$L_r \to 0$ (vertical LM) or $I' \to -\infty$, $dY/dG \to 0$: complete \omterm{def:b1:long-run-model:crowding-out}{crowding out}.
\end{solution}

\begin{solution}{exo:b1:is-lm:refute}
Market interest rates on short-term safe assets falling well below zero as
reserves grew; or each step up in reserves lowering short-term rates as in
normal times.
\end{solution}

\begin{solution}{pb:b1:is-lm:brenland}
\textbf{1.} $c_0 = 225.014 - 0.6(393.264 - 80) = 37.056$.

\textbf{2.} $I_0 = 88.484 + 5 \times 3 = 103.484$.

\textbf{3.} $M/P = 0.5 \times 393.264 - 5 \times 3 = 181.632$.

\textbf{4.} Autonomous spending is $37.056 - 48 + 103.484 + 78.653 + 1.113 = 172.306$, so
$Y = 2.5(172.306 - 5r) = 430.765 - 12.5r$.

\textbf{5.} $r = (0.5Y - 181.632)/5 = 0.1Y - 36.326$.

\textbf{6.} IS: $-1/12.5 = -0.08$; LM: $0.1$.

\textbf{7.} $1/(1 - 0.6) = 2.5$.

\textbf{8.} $1/(0.4 + 5 \times 0.5/5) = 1.111$ for $G$, and $(5/5)/0.9 = 1.111$ for $M/P$.

\textbf{9.} $2.5 \times 10 = 25$ to the left.

\textbf{10.} $Y = 393.264 - 11.111 = 382.153$ and $r = 1.889$.

\textbf{11.} It rises by $5 \times 1.111 = 5.56$, because the interest rate falls.

\textbf{12.} 25: the full Keynesian-cross fall.

\textbf{13.} $11.111/1.111 = 10$.

\textbf{14.} $r = 1.0$ and $I = 103.484 - 5 = 98.484$.

\textbf{15.} $10$.

\textbf{16.} $r = 3$ and $I = 88.484$.

\textbf{17.} Both restore output; the monetary answer leaves a lower rate and
more investment (more capital later), the fiscal one more government
purchases and the same investment.

\textbf{18.} $Y = 359.931$ and $r = -0.333$.

\textbf{19.} $r = 0$ and $Y = 2.5 \times 142.3056 = 355.764$.

\textbf{20.} Nothing: the rate is at its floor and LM is flat there.

\textbf{21.} $2.5$.

\textbf{22.} Until $Y$ reaches $181.632/0.5 = 363.264$, that is by $7.5/2.5 = 3$.

\textbf{23.} $3 + (393.264 - 363.264)/1.111 = 3 + 27 = 30$.

\textbf{24.} Beyond 3, the rate rises off zero and the multiplier falls to
$1.111$: on average $37.5/30 = 1.25$.

\textbf{25.} A fixed marginal propensity to consume out of current income, no
expectations of future taxes or prices, a closed economy, and prices that
do not move.

\textbf{26.} At the zero bound the \omterm{def:b1:is-lm:fiscal-multiplier}{fiscal multiplier} rises from $1.1$ to $2.5$ in the
toy model. That number is the model's: it follows from the assumptions of
answer 25, which the evidence of \cref{ch:b1:stabilisation-policy} does not
support at that size.
\end{solution}
